\section{模型建立}

\subsection{总体技术路线}

数据首先经过输入冻结、字段审计和数据契约检查，再进入 baseline 或候选模型。优化模型的输出经过约束和目标复算；数据分析模型的输出经过正确切分、指标和稳定性检查。混合模型还需要显式记录上下游接口及不确定性传递。

\subsection{统一形式}

在一个抽象的优化子问题中，可写为
\begin{equation}
  \min_{x\in\mathcal{X}}\; f(x;\theta),
  \qquad
  \text{s.t.}\quad g_j(x;\theta)\le 0,\quad j=1,\ldots,m,
  \label{eq:generic-optimization}
\end{equation}
其中 $x$ 为决策变量，$\theta$ 为由题面或数据契约确定的参数，$f$ 为目标函数，$g_j$ 为约束。实际论文必须把抽象符号替换为具体集合、单位、边界和代码变量，并报告是否有全局最优性证据。

对数据分析子问题，应明确估计对象 $Y$、特征/协变量 $X$、切分规则 $S$ 和评价指标 $M$。结果应写成“在指定 $S$ 和 $M$ 下观察到的性能”，而不是脱离数据边界的普遍断言。
